% BAB 2 - LANDASAN TEORI

% Ganti dengan konten dari laporan Word kamu
% Biasanya berisi teori-teori yang relevan dengan penelitian

\subsection{Sistem Informasi}
\indent
Sistem informasi adalah kombinasi dari teknologi, manusia, dan prosedur yang digunakan untuk mengumpulkan, memproses, menyimpan, serta mendistribusikan informasi guna mendukung pengambilan keputusan dan koordinasi dalam organisasi. Menurut Laudon & Laudon dalam Management Information Systems: Managing the Digital Firm (2021), sistem informasi berperan penting dalam meningkatkan efisiensi operasional dan kualitas layanan melalui integrasi data dan proses bisnis. Hal ini sejalan dengan O’Brien dalam Introduction to Information Systems (2004) yang menekankan bahwa sistem informasi modern tidak hanya berfungsi sebagai alat pencatatan, tetapi juga sebagai sarana strategis untuk mendukung keberlangsungan bisnis di era digital. Dalam konteks modern, sistem informasi berbasis web menjadi pilihan utama karena dapat diakses secara fleksibel melalui berbagai perangkat, sehingga mendukung transparansi data dan integrasi proses bisnis [8][9].

\subsection{Poin of Sale (POS)}
\indent
Point of Sale (POS) adalah sistem yang digunakan untuk mencatat transaksi penjualan secara digital dan terintegrasi. POS tidak hanya berfungsi sebagai alat pencatatan, tetapi juga sebagai sistem manajemen yang mampu menghubungkan data transaksi dengan laporan keuangan dan inventori. POS berbasis web memiliki keunggulan dibandingkan POS tradisional karena memungkinkan akses multi device, integrasi data secara real time, serta mempermudah pemilik usaha dalam melakukan monitoring. Penelitian menunjukkan bahwa sistem berbasis web dapat meningkatkan akurasi pencatatan dan efisiensi operasional [1][2]

\subsection{Usaha Laundry}
\indent
Usaha laundry sebagai bagian dari UMKM sering menghadapi kendala dalam pencatatan manual, seperti kesalahan transaksi, keterlambatan informasi, dan laporan yang tidak terstruktur. Digitalisasi layanan laundry melalui sistem POS berbasis web menjadi solusi untuk mengatasi permasalahan tersebut. Penelitian terdahulu menegaskan bahwa sistem laundry berbasis web mampu menyederhanakan proses pencatatan, mempercepat transaksi, serta meningkatkan kepuasan pelanggan melalui antarmuka yang mudah digunakan [3][4][5]. Studi lain juga menekankan bahwa digitalisasi laundry mendukung keberlanjutan usaha kecil dengan menyediakan laporan yang lebih rapi dan terstruktur [6][7].

\subsection{Teori Usability dan User Experience}
\indent
Selain itu, teori usability menurut ISO 9241 11 menekankan bahwa sistem harus memenuhi aspek efektivitas, efisiensi, dan kepuasan pengguna. Prinsip usability sangat relevan dalam pengembangan POS Laundry berbasis web, karena antarmuka yang sederhana dan informatif akan mempermudah karyawan maupun pemilik dalam menggunakan sistem. Dengan demikian, pengembangan sistem POS Laundry berbasis web tidak hanya didukung oleh konsep sistem informasi dan POS, tetapi juga oleh teori usability yang memastikan sistem dapat digunakan secara optimal oleh penggunanya [3].
% Tambahkan subsection lain sesuai kebutuhan
